\section{Related Work}

Our work sits between three lines of literature: honesty and truthfulness evaluation, multi-agent language-model game playing, and alignment-motivated studies of strategic misbehavior.

Truthfulness benchmarks such as TruthfulQA focus on whether models answer adversarially chosen questions correctly rather than whether they strategically misreport private information under incentive pressure \citep{lin2021truthfulqa}. Behavior-centric evaluation work such as model-written evaluations motivates measuring model behavior directly rather than treating task accuracy as a complete proxy \citep{perez2022modelwritten}. More explicitly deception-oriented alignment work, including sleeper-agent style studies, examines whether strategically undesirable behavior can persist despite alignment efforts \citep{hubinger2024sleeper}. Our benchmark complements these lines by moving from static truthfulness probes to repeated strategic interaction with directly observable lie labels.

Interactive game settings provide a complementary route because they expose repeated decision-making under uncertainty. Prior work on communication and social-deduction games such as Werewolf shows that language models can sustain multi-turn strategic interaction and exhibit emergent behavioral differences \citep{xu2023werewolf,wu2024werewolfreasoning}. BS-Bench differs methodologically because Bullshit yields direct lie labels from hidden cards. We do not need to infer deception indirectly from persuasion, voting, or latent-role reasoning.

Recent project-style work has also used bluffing card games such as Coup to study strategic deception and opponent modeling among LLMs \citep{imran2025blinktwise}. That work reinforces the value of deception games as evaluation environments. Bullshit gives us a different empirical target: each turn has a required rank, the actual cards make lie labels directly verifiable, and challenges resolve against hidden-card ground truth. Coup-style settings emphasize role bluffing and opponent modeling with partially latent action semantics; Bullshit gives cleaner turn-level lie and detection labels.

From an evaluation-design perspective, BS-Bench is meant as a benchmark release rather than a prompt-demo anecdote. The contribution is not only the narrative result but also the fixed task definition, fixed prompt conditions, frozen cohort, logged prompt provenance, and tracked artifacts that let later readers replay or reanalyze the same slice. That release orientation matters in strategic settings because small changes in prompts, providers, or inclusion rules can move the observed table equilibrium. It also distinguishes BS-Bench from open-ended social-deduction settings where the label of interest is partly latent and from one-off agent demos where the underlying prompt stack is not fixed. We therefore emphasize direct, replayable behavior labels over broader ecological realism claims.

The benchmark also connects to broader work on oversight, strategic communication, and instruction-following under competitive incentives \citep{irving2018debate,kenton2024weakjudges,ouyang2022instructgpt}. Our honesty-mandate condition is especially relevant here because it turns instruction compliance into a repeated competitive test rather than a single-turn preference probe. The benchmark is therefore narrower than a general theory of deception, but sharper as an environment where incentives, bluff opportunities, and challenge outcomes are all observable and reproducible.
